% =============================================================================
% 03_math_mechanisms.tex — Problem 3: Mathematical Mechanisms (P4 — Triết)
% =============================================================================
\section{Mathematical Mechanisms}
\label{sec:math}

% ---- 3.1 Ridge and Conditioning --------------------------------------------
\subsection{Ridge Regression and Conditioning}
\label{sec:ridge-math}

\begin{definition}[Ridge Objective]
Let $\bX \in \R^{n \times p}$ be the design matrix, $\by \in \R^n$ the
response vector, and $\bbeta \in \R^p$ the coefficient vector.  The Ridge
regression objective is:
\begin{equation}
\label{eq:ridge-obj}
    \hat{\bbeta}^{\mathrm{Ridge}}
    = \argmin_{\bbeta \in \R^p}
      \left\{
          \norm{\by - \bX \bbeta}_2^2
          + \lambda \norm{\bbeta}_2^2
      \right\},
    \qquad \lambda > 0.
\end{equation}
\end{definition}

\paragraph{Derivation.}
Expanding the objective function:
\begin{align}
    L(\bbeta)
    &= (\by - \bX\bbeta)^\top (\by - \bX\bbeta)
       + \lambda\, \bbeta^\top \bbeta \notag \\
    &= \by^\top\by
       - 2\, \bbeta^\top \bX^\top \by
       + \bbeta^\top \bX^\top \bX\, \bbeta
       + \lambda\, \bbeta^\top \bbeta.
    \label{eq:ridge-expand}
\end{align}

Taking the gradient with respect to $\bbeta$ and setting it to zero:
\begin{align}
    \frac{\partial L}{\partial \bbeta}
    &= -2\, \bX^\top \by + 2\, \bX^\top \bX\, \bbeta
       + 2\lambda\, \bbeta
    = \bm{0} \notag \\
    \Longrightarrow \quad
    (\bX^\top \bX + \lambda\, \bm{I}_p)\, \bbeta
    &= \bX^\top \by.
    \label{eq:ridge-normal}
\end{align}

Since $\bX^\top \bX$ is positive semi-definite and $\lambda > 0$, the matrix
$\bX^\top \bX + \lambda\, \bm{I}_p$ is \textbf{strictly positive definite},
hence always invertible.  The unique closed-form solution is:
\begin{equation}
\label{eq:ridge-solution}
    \boxed{%
        \hat{\bbeta}^{\mathrm{Ridge}}
        = \bigl(\bX^\top \bX + \lambda\, \bm{I}_p\bigr)^{-1} \bX^\top \by.
    }
\end{equation}

\begin{remark}[Conditioning]
If $\sigma_1 \geq \sigma_2 \geq \cdots \geq \sigma_p > 0$ are the singular
values of $\bX$, then the condition number of $\bX^\top\bX$ is
$\kappa(\bX^\top\bX) = (\sigma_1 / \sigma_p)^2$.
Adding the ridge penalty changes the eigenvalues to
$\sigma_j^2 + \lambda$, so:
\[
    \kappa\bigl(\bX^\top\bX + \lambda\bm{I}\bigr)
    = \frac{\sigma_1^2 + \lambda}{\sigma_p^2 + \lambda}
    \;\leq\; \frac{\sigma_1^2}{\sigma_p^2}
    = \kappa(\bX^\top\bX).
\]
As $\lambda$ increases, the condition number decreases monotonically toward~$1$,
stabilising the inversion.
\end{remark}

% TODO: Connect this theory to the numerical condition numbers from
%       Table~\ref{tab:condition-numbers}. Show that the computed ratio
%       matches the formula above for at least one lambda value.

% ---- 3.2 Lasso Optimality and Soft Thresholding -----------------------------
\subsection{Lasso Optimality and Soft Thresholding}
\label{sec:lasso-math}

\begin{definition}[Lasso Objective]
\begin{equation}
\label{eq:lasso-obj}
    \hat{\bbeta}^{\mathrm{Lasso}}
    = \argmin_{\bbeta \in \R^p}
      \left\{
          \frac{1}{2n} \norm{\by - \bX \bbeta}_2^2
          + \lambda \norm{\bbeta}_1
      \right\}.
\end{equation}
\end{definition}

\paragraph{Subgradient optimality conditions.}
Because $\norm{\bbeta}_1 = \sum_{j=1}^{p} \abs{\beta_j}$ is not differentiable
at $\beta_j = 0$, we use subgradient calculus.  The subgradient optimality
condition for coordinate~$j$ is:
\begin{equation}
\label{eq:lasso-subgradient}
    -\frac{1}{n} \bx_j^\top \bigl(\by - \bX \hat{\bbeta}\bigr)
    + \lambda\, \partial\abs{\hat{\beta}_j} \ni 0,
\end{equation}
where $\partial\abs{\hat\beta_j}$ is the subdifferential of the absolute value:
\[
    \partial\abs{\hat\beta_j} =
    \begin{cases}
        \{+1\}        & \text{if } \hat\beta_j > 0, \\
        \{-1\}        & \text{if } \hat\beta_j < 0, \\
        [-1, +1]      & \text{if } \hat\beta_j = 0.
    \end{cases}
\]

\paragraph{Special case: orthonormal design.}
When $\frac{1}{n}\bX^\top\bX = \bm{I}_p$, each coordinate decouples and the
Lasso admits the closed-form \textbf{soft-thresholding} solution:
\begin{equation}
\label{eq:soft-threshold}
    \boxed{%
        \hat\beta_j^{\mathrm{Lasso}}
        = \mathrm{sign}\!\bigl(\hat\beta_j^{\mathrm{OLS}}\bigr)
          \cdot \max\!\bigl(\abs{\hat\beta_j^{\mathrm{OLS}}} - \lambda,\; 0\bigr).
    }
\end{equation}

This means:
\begin{itemize}
    \item If $\abs{\hat\beta_j^{\mathrm{OLS}}} > \lambda$: the coefficient is
          \emph{shrunk} toward zero by exactly $\lambda$.
    \item If $\abs{\hat\beta_j^{\mathrm{OLS}}} \leq \lambda$: the coefficient
          is set \emph{exactly to zero} — the variable is dropped.
\end{itemize}

% TODO: Add a brief interpretation of why the L1 geometry (diamond constraint)
%       causes exact zeros while L2 (sphere constraint) does not.

% ---- 3.3 Connecting Algebra to Evidence -------------------------------------
\subsection{Connecting Algebra to Evidence}
\label{sec:algebra-evidence}

% TODO: Pick ONE specific numerical result from Section~\ref{sec:ols-ridge-lasso}
%       and show how it follows from the theory above.
%
% Example ideas:
%   (a) Take the condition number from Table~\ref{tab:condition-numbers} at
%       lambda_min. Verify: kappa(X'X + lambda*I) ≈ (sigma_1^2 + lambda)/(sigma_p^2 + lambda).
%   (b) Pick a variable zeroed out by Lasso. Show its OLS coefficient magnitude
%       is < lambda, consistent with the soft-thresholding rule.
%   (c) Compare Ridge and Lasso coefficients for a specific variable and
%       explain the difference using the L1 vs L2 penalty geometry.

\textit{TODO: Provide a concrete numerical example linking the theory from
Sections~\ref{sec:ridge-math}--\ref{sec:lasso-math} to the empirical results
from Section~\ref{sec:ols-ridge-lasso}.}
